\chap{Abstract} 


%%%% your abstract goes here (word limit: 350)
Access to scientific discovery sustains trust in scientific institutions and the public's willingness to continue investing in them. But accessibility is not a fixed target; what makes a paper accessible depends on who is reading it. A parent deciding whether to vaccinate their child and a pediatrician advising that family might read the same clinical study. However, the parent needs the bottom-line risk explained in plain terms while the pediatrician needs the contraindications and dosing details. A summary written for one fails the other.
When scientific communication fails, trust in the science itself erodes. Current summarization methods largely sidestep this problem by assuming a generic reader: a fixed reading level, a fixed level of detail, a fixed notion of what counts as important. \textbf{My research reduces barriers to accessing scientific knowledge by developing methods to generate and evaluate summarization systems that serve the informational and practical needs of the end user. }


If accessibility of science is user-dependent, then both the systems that produce summaries and the methods that judge them must be user-dependent too. The first part of my thesis focuses on generating summaries that adapt to users’ needs. I introduce methods for controllable summarization that let users specify the structure and technicality of a summary. I show that these methods outperform strong baselines and substantially improve the performance of smaller, low-compute language models. The second part focuses on evaluating whether those summaries actually meet users' needs. I show that the current standard of evaluation has low agreement with human judgments on both readability and information satisfaction. I then develop LLM-based evaluation methods that capture deeper attributes of readability such as explanation of technical concepts and propose methods to improve LLM-as-judges in a low data setting. Together, these contributions move scientific summarization away from one-size-fits-all systems, toward tools designed around the users they are meant to serve, reducing the gap between scientific knowledge being available and being genuinely accessible.




%% list of keywords seperated by comma
\keywords{summarization, natural language processing, human-computer interaction, evaluation, natural language generation, language models, artificial intelligence}


%%%%  committee members (add it right after the abstract w/o page break)
\begin{singlespace}

    %% if you have co-advisor, add here w/ \vspace{0.1in} as shown below
    %% alternatively you can use minipage environment to put side-by-side
    \section*{Primary reader and thesis advisor}
    
    Dr. Mark Dredze \\
    Professor\\
    Department of Computer Science\\
    Johns Hopkins University, Baltimore MD 


    \section*{Secondary readers}
    
    Dr. Daniel Khashabi\\
    Assistant Professor\\
    Department of Computer Science \\
    Johns Hopkins University, Baltimore, MD 
    
    \vspace{0.1in}
    
    Dr. Sharon Levy \\
    Assistant Professor\\
    Department of Computer Science \\
    Rutgers University, New Brunswick, NJ

    %% you can add more readers if you have them on your committee 
    %% use \vspace{0.1in} in between members for clarity
    %% you can also place committee members side-by-side using `minipage`


\end{singlespace}